To perform the bottleneck analysis, we utilized nvprof to perform a kernel-level profiling of our program. What we wanted to discover was what portions or kernel calls 
were responsible the largest portion of program time. \strongtabref{tab:profile_config} shows the configuration for the model and training set size used in the profiling. The batch size
parameter is only applicable to the second profiling result, after optimizations were introduced.
\begin{table}[H]
\centering
\small
\renewcommand{\arraystretch}{1.2}
\setlength{\tabcolsep}{6pt}
\begin{tabular}{lc}
\hline
\textbf{Parameter} & \textbf{Value} \\
\hline
Names               & 10{,}000 \\
Batch Size          & 512 \\
Number of Layers    & 8 \\
Embedding Dimension & 512 \\
Block Size          & 64 \\
Number of Heads     & 16 \\
\hline
\end{tabular}
\caption{Profiling Configuration Parameters}
\label{tab:profile_config}
\end{table}

\begin{table}[t]
\centering
\scriptsize
\renewcommand{\arraystretch}{1.1}
\setlength{\tabcolsep}{4pt}
\begin{tabular}{lccccc}
\hline
\textbf{Kernel} & \makecell{\textbf{Time} \\ \textbf{(\%)}} & \makecell{\textbf{Total} \\ \textbf{(s)}} & \textbf{Calls} & \makecell{\textbf{Avg} \\ \textbf{($\mu$s)}} & \makecell{\textbf{Max} \\ \textbf{($\mu$s)}} \\
\hline
linear & 80.86 & 70.32 & 490k & 143.51 & 464.80 \\
RMSNorm & 10.33 & 8.99 & 170k & 52.86 & 82.27 \\
self-attn & 7.33 & 6.38 & 80k & 79.70 & 162.91 \\
cross-ent. & 0.85 & 0.74 & 10k & 73.68 & 149.54 \\
resid. add & 0.25 & 0.22 & 160k & 1.36 & 1.95 \\
ReLU & 0.13 & 0.11 & 80k & 1.37 & 2.14 \\
embed. look. & 0.02 & 0.02 & 10k & 1.61 & 2.18 \\
\hline
\end{tabular}
\caption{CUDA Kernel Profiling Results for Parallel Implementation Pre-Optimizations}
\label{tab:profile_pre_optimization}
\end{table}

What we noticed in \strongtabref{tab:profile_pre_optimization} were two things:
\begin{enumerate}
    \item Over $80\%$ of execution time was spent in the \texttt{linear\_kernel}.
    \item The average time per kernel was reasonable, but we had many invocations of them.
\end{enumerate}
From these investigations, we concluded that the most effective path forward was to adopt optimization strategies standard with modern large language model (LLM) training practices. Our first focus was improving the performance of the linear layer, which is dominated by matrix multiplication. By implementing a tiling strategy, as introduced in lecture, and transitioning from double-precision to single-precision (float) arithmetic, we reduced the average runtime of the \texttt{linear\_kernel} by approximately $81.87\times$.

Building on this improvement, we next targeted batching, a cornerstone of modern training systems. Processing individual tokens or sequences underutilizes the parallel compute capabilities of GPUs. By restructuring our implementation to process multiple sequences simultaneously in a single batch, we were able to more effectively leverage GPU parallelism and further improve overall throughput.

\strongfigref{result:ablation} presents the ablation study used to evaluate the impact of our optimizations. Overall, the results demonstrate that these changes effectively reduce the runtime of the forward pass.
One notable observation is that the tiling optimization, in isolation, led to an increase in runtime. This behavior is expected for smaller problem sizes where batching is not applied. In such cases, the overhead associated with thread synchronization and shared memory transfers outweighs the benefits of improved data locality, resulting in a net slowdown.

\begin{figure}[t]
    \centering
    \includegraphics[width=\linewidth]{cuda_ablation_names.png}
    \caption{Ablation Study of Implemented Optimizations}
    \label{result:ablation}
\end{figure}

Our profiling results for the optimized implementation can be found in \strongappref{appendix:profile}.
